\section{Machine Learning Algorithms}

The Python part of this exercise is available in \myhref{https://github.com/LosDelDGIIM/LosDelDGIIM.github.io/blob/main/subjects/Erasmus-DUE/GKI/Exercises/Exercise03.ipynb}{this Jupyter Notebook}.

\subsection*{The Intelligent Spam Filter}

Imagine you have trained a model to distinguish important emails from annoying advertisements. To evaluate the quality of your filter, you ran a test on 100 emails and obtained the results shown in Table~\ref{tab:confusion-matrix-ejs}.
\begin{table}
    \centering
    \caption{Confusion matrix for the spam filter model.}
    \begin{tabular}{l||c|c}
        & \textbf{Actual: Spam} & \textbf{Actual: Not Spam} \\
        \hline \hline
        \textbf{Model says: Spam} & 30 (TP) & 10 (FP) \\
        \textbf{Model says: Not Spam} & 20 (FN) & 40 (TN) \\
    \end{tabular}
    \label{tab:confusion-matrix-ejs}
\end{table}

\subsubsection*{Confusion Matrix}
\begin{ejercicio}
    State the formulas for Accuracy, Precision, Recall, and F1-Score using $TP$, $FP$, $TN$, $FN$, and specify their respective ranges.
    \begin{itemize}
        \item \ul{Accuracy}: It is ranged in $[0, 1]$.
        \begin{equation*}
            \text{Accuracy} = \frac{TP + TN}{TP + FP + TN + FN}
        \end{equation*}
        \item \ul{Precision}: It is ranged in $[0, 1]$.
        \begin{equation*}
            \text{Precision} = \frac{TP}{TP + FP}
        \end{equation*}
        \item \ul{Recall}: It is ranged in $[0, 1]$.
        \begin{equation*}
            \text{Recall} = \frac{TP}{TP + FN}
        \end{equation*}
        \item \ul{F1-Score}: It is ranged in $[0, 1]$.
        \begin{align*}
            \text{F1-Score} &= 2 \cdot \frac{\text{Precision} \cdot \text{Recall}}{\text{Precision} + \text{Recall}}
        \end{align*}
    \end{itemize}
\end{ejercicio}
\begin{ejercicio}
    Calculate the values using the table above. What do these values specifically say about the reliability of your filter in everyday use?
    \begin{itemize}
        \item \ul{Accuracy}: $\dfrac{30 + 40}{30 + 10 + 40 + 20} = \dfrac{70}{100} = 0.7$.
        \item \ul{Precision}: $\dfrac{30}{30 + 10} = \dfrac{30}{40} = 0.75$.
        \item \ul{Recall}: $\dfrac{30}{30 + 20} = \dfrac{30}{50} = 0.6$.
        \item \ul{F1-Score}: $2 \cdot \dfrac{0.75 \cdot 0.6}{0.75 + 0.6} = 2 \cdot \dfrac{0.45}{1.35} = 2 \cdot \dfrac{1}{3} = \dfrac{2}{3} \approx 0.67$.
    \end{itemize}

    These values indicate that while the filter is reasonably accurate (70\%), it has a moderate precision (75\%) and a lower recall (60\%). This means that while most of the emails it identifies as spam are indeed spam, it misses a significant portion (40\%) of actual spam emails, which could lead to important emails being overlooked. In addition, the F1-Score of approximately 0.67 suggests that there is room for improvement in balancing precision and recall for better overall performance in everyday use.
\end{ejercicio}
\begin{ejercicio}
    To get a complete picture, we will also use the Fowlkes-Mallows Index (FMI). What is the formula, and calculate the value.
    \begin{equation*}
        \text{FMI} = \sqrt{\text{Precision} \cdot \text{Recall}} = \sqrt{0.75 \cdot 0.6} \approx 0.67
    \end{equation*}
\end{ejercicio}
\begin{ejercicio}
    Which value is more sensitive to the missed spam emails (FN) in your example?

    The Recall is more sensitive to the missed spam emails (FN) because it directly measures the proportion of actual spam emails that were correctly identified by the model. A high number of false negatives (missed spam emails) will significantly decrease the Recall, indicating that the model is not effectively capturing all the spam emails in the dataset.
\end{ejercicio}
\begin{ejercicio}
    Suppose your filter has an accuracy of 100\% on the training data, but only 70\% in the test table above. What phenomenon is at play here? Briefly explain why memorized emails hinder generalization.

    The phenomenon at play here is called \emph{overfitting}. Overfitting occurs when a model learns the training data too well, including its noise and outliers, to the point that it performs exceptionally on the training data but poorly on unseen test data. Memorized emails hinder generalization because the model has essentially ``memorized'' specific examples from the training set rather than learning the underlying patterns that can be applied to new, unseen data. As a result, when the model encounters new emails in the test set, it fails to generalize and correctly classify them, leading to a significant drop in accuracy.
\end{ejercicio}
\begin{ejercicio}
    Explain the variance-bias dilemma: Which of these two factors predominates when the model begins to learn the noise in the data (overfitting), and which predominates when the model is oversimplified (underfitting)?

    The variance-bias dilemma refers to the trade-off between two sources of error in machine learning models: bias and variance.
    \begin{itemize}
        \item \ul{Bias}: Bias is the error introduced by approximating a real-world problem, which may be complex, by a simplified model. High bias can lead to underfitting, where the model is too simple to capture the underlying patterns in the data.
        \item \ul{Variance}: Variance is the error introduced by the model's sensitivity to small fluctuations in the training data. High variance can lead to overfitting, where the model learns not only the underlying patterns but also the noise in the training data.
    \end{itemize}
\end{ejercicio}
\begin{ejercicio}
    How does the complexity of the model change when you try to reduce the bias?

    When you try to reduce the bias, you typically increase the complexity of the model. This is because reducing bias often involves adding more parameters or using a more flexible model that can capture more complex patterns in the data. However, increasing the complexity of the model can also lead to higher variance, which may result in overfitting if not managed properly.
\end{ejercicio}

\subsubsection*{Data Imbalance (Resampling)}
\begin{ejercicio}
    The needle-in-a-haystack problem: Your dataset contains 9900 normal emails and only 100 spam emails. Explain why a model without resampling tends to simply always predict Not Spam. What would the recall value be in this extreme case?

    In this extreme case, a model without resampling would likely learn to predict Not Spam for all emails because it would achieve a high accuracy of 99\% by doing so (9900 out of 10000 emails are Not Spam). However, this approach completely ignores the minority class (Spam), leading to a recall value of 0\% for the Spam class. This is because the model would fail to identify any of the actual spam emails, resulting in zero true positives (TP) and thus a recall of:
    \begin{equation*}
        \text{Recall} = \frac{TP}{TP + FN} = \frac{0}{0 + 100} = 0
    \end{equation*}

    Resampling techniques are crucial to address this imbalance and allow the model to learn to identify spam emails effectively. Some options include oversampling the minority class or undersampling the majority class to create a more balanced dataset for training.
\end{ejercicio}
\begin{ejercicio}
    Why is it crucial to apply resampling to the training data and not to the test data? What would happen to the validity of your metrics if the test table were artificially balanced?

    It is crucial to apply resampling only to the training data because the test data should reflect the real-world distribution of classes to provide an accurate evaluation of the model's performance. If the test table were artificially balanced, it would not represent the actual conditions under which the model will operate, leading to misleading metrics. For instance, if the test set is balanced, the model might appear to perform well in terms of recall and precision for both classes, but in reality, it may struggle to identify spam emails in a real-world scenario where they are much less frequent. This would compromise the validity of the metrics and give a false sense of confidence in the model's performance.
\end{ejercicio}

\subsection*{Automatic Fruit Sorting (kNN)}

You are using a sensor that classifies fruit based on two features: weight ($x_1$) and hardness ($x_2$).

\subsubsection*{The k-Nearest-Neighbors Algorithm}
\begin{ejercicio}
    An unknown piece of fruit lands on the conveyor belt. Briefly describe the three algorithmic steps that kNN now performs to determine whether it is an apple or a pear.
    \begin{itemize}
        \item \ul{Step 1: Calculate Distances}: The algorithm calculates the distance between the unknown fruit and all the known fruits in the training dataset using a distance metric.
        \item \ul{Step 2: Identify Neighbors}: The algorithm identifies the $k$ nearest neighbors (the $k$ closest known fruits) based on the calculated distances.
        \item \ul{Step 3: Make a Prediction}: The algorithm makes a prediction for the unknown fruit by taking a majority vote among the $k$ nearest neighbors. The class (apple or pear) that appears most frequently among the neighbors is assigned to the unknown fruit.
    \end{itemize}
\end{ejercicio}
\begin{ejercicio}
    What are the formulas for the Euclidean distance and the Manhattan distance between two points $P(x_1, x_2)$ and $Q(y_1, y_2)$?

    The formulas for the distances are as follows:
    \begin{itemize}
        \item \ul{Euclidean Distance}:
        \begin{equation*}
            d_{\text{Euclidean}}(P, Q) = \sqrt{(x_1 - y_1)^2 + (x_2 - y_2)^2}
        \end{equation*}
        \item \ul{Manhattan Distance}:
        \begin{equation*}
            d_{\text{Manhattan}}(P, Q) = |x_1 - y_1| + |x_2 - y_2|
        \end{equation*}
    \end{itemize}
\end{ejercicio}
\begin{ejercicio}
    The new fruit has the coordinates $(5, 2)$. A known reference apple is located at $(10, 5)$. Calculate the distance using both of the above-mentioned distance methods.
    \begin{itemize}
        \item \ul{Euclidean Distance}:
        \begin{equation*}
            d_{\text{Euclidean}}((5, 2), (10, 5)) = \sqrt{(5 - 10)^2 + (2 - 5)^2} = \sqrt{25 + 9} = \sqrt{34} \approx 5.83
        \end{equation*}
        \item \ul{Manhattan Distance}:
        \begin{equation*}
            d_{\text{Manhattan}}((5, 2), (10, 5)) = |5 - 10| + |2 - 5| = 5 + 3 = 8
        \end{equation*}
    \end{itemize}
\end{ejercicio}

\subsection*{k-Nearest Neighbors}

\begin{ejercicio}
    What are the fundamental differences in the operating principles of kNN, k-Means, and k-Medians, and what types of tasks are they typically used for (classification, regression, clustering)?
    \begin{description}
        \item \ul{kNN (k-Nearest Neighbors)}: kNN is a supervised learning algorithm used for classification and regression tasks. It classifies a new data point based on the majority class among its $k$ nearest neighbors in the feature space. For regression, it predicts the value based on the average of the target values of the $k$ nearest neighbors.

        \item \ul{k-Means}: k-Means is an unsupervised learning algorithm used for clustering tasks. It partitions the data into $k$ clusters by assigning each data point to the cluster with the nearest mean (centroid). The algorithm iteratively updates the centroids until convergence. k-Means is sensitive to outliers because the mean can be significantly affected by extreme values.

        \item \ul{k-Medians}: k-Medians is also an unsupervised learning algorithm used for clustering tasks, similar to k-Means. However, instead of using the mean to represent each cluster, it uses the median, which makes it more robust to outliers compared to k-Means.
    \end{description}
\end{ejercicio}
\begin{ejercicio}
    How exactly does kNN work step by step?
    \begin{enumerate}
        \item \ul{Distance Calculation}: For a given new data point, the algorithm calculates the distance between this point and all the points in the training dataset using a chosen distance metric (e.g., Euclidean, Manhattan).
        \item \ul{Neighbor Identification}: The algorithm identifies the $k$ nearest neighbors to the new data point based on the calculated distances.
        \item \ul{Prediction}: The algorithm makes a prediction for the new data point depending on the task:
        \begin{itemize}
            \item For classification tasks, it assigns the class that is most common among the $k$ nearest neighbors.
            \item For regression tasks, it calculates the average (or median) of the target values of the $k$ nearest neighbors and assigns this value to the new data point.
        \end{itemize}
    \end{enumerate}
\end{ejercicio}
\begin{ejercicio}
    What distance metrics do you know, and provide their mathematical forms?
    \begin{itemize}
        \item \ul{Euclidean Distance}:
        \begin{equation*}
            d_{\text{Euclidean}}(P, Q) = \sqrt{\sum_{i=1}^{n} (x_i - y_i)^2}
        \end{equation*}
        \item \ul{Manhattan Distance}:
        \begin{equation*}
            d_{\text{Manhattan}}(P, Q) = \sum_{i=1}^{n} |x_i - y_i|
        \end{equation*}
        \item \ul{Minkowski Distance}:
        \begin{equation*}
            d_{\text{Minkowski}}(P, Q) = \sqrt[p]{\sum_{i=1}^{n} |x_i - y_i|^p }
        \end{equation*}
        where $p$ is a parameter that can be adjusted to obtain different distance metrics (e.g., $p=1$ for Manhattan, $p=2$ for Euclidean).
    \end{itemize}
\end{ejercicio}
\begin{ejercicio}
    What are the advantages and disadvantages of kNN, k-Means, and k-Medians in terms of robustness to outliers, computational effort, and interpretability of results?
    \begin{description}
        \item \ul{kNN}:
        \begin{itemize}
            \item Advantages: Simple to understand and implement, can capture complex relationships in the data, and can be used for both classification and regression tasks.
            \item Disadvantages: Computationally expensive for large datasets, sensitive to irrelevant features and the choice of $k$, and does not provide a model or insights into feature importance.
        \end{itemize}

        \item \ul{k-Means}:
        \begin{itemize}
            \item Advantages: Efficient for large datasets, easy to implement, and provides clear cluster assignments.
            \item Disadvantages: Sensitive to outliers, requires the number of clusters ($k$) to be specified in advance, and may converge to local minima.
        \end{itemize}

        \item \ul{k-Medians}:
        \begin{itemize}
            \item Advantages: More robust to outliers than k-Means, as it uses the median instead of the mean to represent clusters.
            \item Disadvantages: Computationally more expensive than k-Means due to the need to calculate medians, and also requires the number of clusters ($k$) to be specified in advance.
        \end{itemize}
    \end{description}
\end{ejercicio}
\begin{ejercicio}
    How do you choose the optimal value for ``k'' in kNN, k-Means, and k-Medians, and what methods are there to evaluate the performance of the models?
    \begin{itemize}
        \item \ul{kNN}: The optimal value of $k$ can be chosen using techniques such as cross-validation, where the dataset is split into training and validation sets to evaluate the performance of the model for different values of $k$. The value of $k$ that yields the best performance (e.g., highest accuracy, lowest error) on the validation set is selected.

        \item \ul{k-Means and k-Medians}: The optimal number of clusters ($k$) can be determined using methods such as the Elbow Method, which involves plotting the within-cluster sum of squares (WCSS) against different values of $k$ and looking for an ``elbow'' point where the rate of decrease sharply changes.
    \end{itemize}
\end{ejercicio}
\begin{ejercicio}
    In what scenarios would k-Medians be preferable to k-Means, and what specific challenges arise in the implementation and application of k-Medians?
    
    k-Medians would be preferable to k-Means in scenarios where the dataset contains outliers or is not normally distributed, as k-Medians is more robust to outliers due to its use of the median instead of the mean. However, challenges in implementing k-Medians include increased computational complexity compared to k-Means, as calculating the median can be more time-consuming than calculating the mean (since it may require sorting the data). Additionally, like k-Means, k-Medians also requires the number of clusters ($k$) to be specified in advance, which can be difficult to determine without prior knowledge of the data.
\end{ejercicio}